\begin{question}[index=17,points=10]
设函数 $y(x)$ 是微分方程 $\displaystyle y'+\frac{1}{2\sqrt x}y=\frac{\sqrt x}{2}+1$
满足 $y(1)=2$ 的解，求曲线 $y=y(x)$ 的渐近线。
\end{question}
\QuestionBreak\WorkSpace

\begin{question}[index=18,points=12]
设 $f(x)$ 是可导函数，且 $f(0)=0$，$\displaystyle g(x)=\int_0^1 xf(tx)\dd t$，
并满足方程 $f'(x)+g'(x)=x$。求由曲线 $y=f(x)$，$y=\mathrm e^{-x}$
及直线 $x=0$，$x=2$ 所围成的平面图形的面积 $S$。
\end{question}
\QuestionBreak\WorkSpace\WorkPageBreak

\begin{question}[index=19,points=12]
设 $f(x)$ 在 $[0,1]$ 上具有二阶导数，$f(0)=f(1)=0$，$f''(x)<0$，$0\le f(x)\le1$。证明：
\begin{enumerate}[label=（\arabic*）,leftmargin=2.5em,labelsep=.4em,itemsep=3pt,topsep=3pt]
\item 存在 $\xi\in(0,1)$，使得对任意 $x\in(0,\xi)$，有 $f'(x)>0$；
\item $\displaystyle\int_0^1\sqrt{1+[f'(x)]^2}\dd x<3$。
\end{enumerate}
\end{question}
\QuestionBreak\WorkSpace

\begin{question}[index=20,points=12]
求 $\displaystyle I=\int_C[\mathrm e^x\sin y-b(x+y)]\dd x+(\mathrm e^x\cos y-ax)\dd y$，
其中 $a,b$ 为正常数，$C$ 为从点 $A(0,2a)$ 沿曲线 $x=-\sqrt{2ay-y^2}$ 到点 $O(0,0)$ 的有向弧。
\end{question}
\QuestionBreak\WorkSpace\WorkPageBreak

\begin{question}[index=21,points=12]
已知 $1,1,-1$ 是 $3$ 阶实对称矩阵 $\symbfit A$ 的 $3$ 个特征值，
$\symbfit\eta_1=(1,1,1)^{\mathrm T}$，$\symbfit\eta_2=(2,2,1)^{\mathrm T}$
是 $\symbfit A$ 的属于特征值 $\lambda_1=\lambda_2=1$ 的特征向量。
\begin{enumerate}[label=（\arabic*）,leftmargin=2.5em,labelsep=.4em,itemsep=3pt,topsep=3pt]
\item 求 $\symbfit A$ 的属于特征值 $\lambda_3=-1$ 的特征向量；
\item 求可逆矩阵 $\symbfit P$，使得
$(\symbfit A\symbfit P+\symbfit P)^{\mathrm T}(\symbfit A\symbfit P+\symbfit P)$ 为对角矩阵。
\end{enumerate}
\end{question}
\QuestionBreak\WorkSpace

\begin{question}[index=22,points=12]
设总体 $X$ 的概率密度为
\[
f(x;\theta)=\begin{cases}
1,&\theta-\dfrac12\le x\le\theta+\dfrac12,\\[3pt]
0,&\text{其他},
\end{cases}
\]
其中 $-\infty<\theta<+\infty$，$X_1,X_2,\ldots,X_n$ 为取自总体 $X$ 的简单随机样本，并记
\[
X_{(1)}=\min\{X_1,X_2,\ldots,X_n\},\qquad X_{(n)}=\max\{X_1,X_2,\ldots,X_n\}.
\]
\begin{enumerate}[label=（\arabic*）,leftmargin=2.5em,labelsep=.4em,itemsep=3pt,topsep=3pt]
\item 求参数 $\theta$ 的矩估计量 $\hat\theta_M$ 和全部最大似然估计量 $\hat\theta_L$；
\item 判断 $\displaystyle\frac{X_{(1)}+X_{(n)}}2$ 是否为 $\theta$ 的无偏估计量，并说明理由。
\end{enumerate}
\end{question}
\WorkSpace
